\NativeSlot{40}{36}{%
% p042_s01

\ExamQuestion{7.}{5}%
商店购进6箱火龙果，每箱24千克，如果分装成每袋8千克，一共可以分装\mbox{（\hspace{1.8em}）}袋；把这些火龙果按每袋12元售卖，全部售出一共收入\mbox{（\hspace{1.8em}）}元。\par
}
\NativeSlot{90}{42}{%
% p042_s02

\ExamQuestion{8.}{5}%
右图中，有\mbox{（\hspace{1.8em}）}个三角形，\mbox{（\hspace{1.8em}）}个平行四边形，\mbox{（\hspace{1.8em}）}个梯形。\par
\BookFigure{figures/p042_s02-01.tex}
}
\NativeSlot{166}{36}{%
\ExamGroup{二、选择（将正确答案的序号填在括号里）。{\normalfont\kaishu\fontsize{10.8bp}{14bp}\selectfont （每题2分，共10分）}}

% p042_s03

\ExamQuestion{1.}{5}%
与\(720\div(8\times9)\)结果相等的算式是\mbox{（\hspace{1.8em}）}。\par
\ExamOptions{A. \(720\div8\times9\)}{B. \(720\div8\div9\)}{C. \(720\times8\div9\)}{D. \(720\div9\times8\)}
}
\NativeSlot{216.0}{36}{%
% p042_s04

\ExamQuestion{2.}{5}%
\(25\times(a+4)\)与\(25\times a+25\)的结果相差\mbox{（\hspace{1.8em}）}。\par
\ExamOptions{A. \(24a\)}{B. \(a\)}{C. 4}{D. 75}
}
\NativeSlot{265.0}{43}{%
% p042_s05

\ExamQuestion{3.}{5}%
在同一平面内，如果直线\(a\)垂直于直线\(b\)，直线\(b\)垂直于直线\(c\)，那么直线\(a\)和\(c\)\mbox{（\hspace{1.8em}）}。\par
\ExamOptions{A. 互相平行}{B. 互相垂直}{C. 相交但不垂直}{D. 无法确定}
}
\NativeSlot{340}{36}{%
% p042_s06

\ExamQuestion{4.}{5}%
下面每组中的三根小棒，能围成三角形的是\mbox{（\hspace{1.8em}）}（单位：厘米）。\par
\ExamOptions{A. 4、5、9}{B. 3、4、8}{C. 5、6、10}{D. 2、2、5}
}
\NativeSlot{390}{36}{%
% p042_s07

\ExamQuestion{5.}{5}%
一个三角形的最小角是\(50^\circ\)，这个三角形一定是\mbox{（\hspace{1.8em}）}。\par
\ExamOptions{A. 锐角三角形}{B. 直角三角形}{C. 钝角三角形}{D. 无法判断}
}
\NativeSlot{441.0}{59}{%
\ExamGroup{三、计算。{\normalfont\kaishu\fontsize{10.8bp}{14bp}\selectfont （共24分）}}

% p042_s08

\ExamQuestion{1.}{5}%
直接写得数。（8分）\par
\noindent \NativeOralCell{25\times3\times4=}\NativeOralCell{125\times8=}\NativeOralCell{23\times10=}\NativeOralCell{200-100\div10=}\par
\noindent \NativeOralCell{360\div60=}\NativeOralCell{11\times20=}\NativeOralCell{(12+8)\times5=}\NativeOralCell{80\div(8\times2)=}\par
}
\NativeSlot{541.0}{71}{%
\noindent \textbf{2.}\quad 计算下面各题，能简算的要简算。（16分）\par\smallskip
\par\noindent%
\ExamTile{.485}{%
% p042_s09

\noindent \(640\div[16\times(19-14)]\)\par


}\hfill%
\ExamTile{.485}{%
% p042_s10

\noindent \(25\times32\times25\)\par


}\par
}
\NativeSlot{667}{68}{%
\par\noindent%
\ExamTile{.485}{%
% p042_s11

\noindent \(99\times47+47\)\par


}\hfill%
\ExamTile{.485}{%
% p042_s12

\noindent \(35\times98\)\par


}\par
}
